\section{Johdanto}

Merkintätapoja.
\begin{itemize}
    \item $\approx$ homeomorfismi
    \item $\linkki,\solmu$ linkki ja solmu
    \item $N(\linkki),N(\solmu)$ linkin ja solmun putkiympäristö
    \item $M_i,L_i$ linkin osasen pituus- ja leveyspiiri
    \item $\#$ topologisten monistojen yhtenäinen summa
    \item 
\end{itemize}

Tämän tekstin lähteenä oleva tiedosto on opinnäytteen pohja, jota voi käyttää
kandidaatintyössä, diplomityössä ja lisensiaatintyössä. Tekstin lähteenä oleva
tiedosto on kirjoitettu \LaTeX-tiedoston rakenteen opiskelemista ajatellen.
Tiedoston kommentit sisältävät tietoa, joka on hyödyllistä opinnäytettä
kirjoitettaessa. 

%% Esimerkki luettelosta. Lyhyt ajatusviiva on käytössä luettelossa, ja se on
%% pituudeltaan en dash. Merkitään latex-koodissa --. 
Johdanto selvittää samat asiat kuin tiivistelmä, mutta laveammin. Johdannossa
kerrotaan yleensä seuraavat asiat

\begin{itemize}
\item[--]Tutkimuksen taustaa ja tutkimusaiheen yleisluonteinen esittely
\item[--]Tutkimuksen tavoitteet
\item[--]Pääkysymys ja osaongelmat
\item[--]Tutkimuksen rajaus ja keskeiset käsitteet.
\end{itemize}

Lyhyiden opinnäytteiden johdannot ovat yleensä liian pitkiä, joten johdannon
paisuttamista on vältettävä. Diplomityöhön sopii johdanto, joka on 2--4 sivua.
%% tässä on myös lyhyt ajatusviiva l. en dash.
Kandidaatintyön johdannon on oltava diplomityön johdantoa lyhyempi. Sopivasti
tiivistetty johdanto ei kaipaa alaotsikoita.

\begin{lause}
    \label{lause:päätulos}
    (i) Pelkistetty algebrallinen käyrä $V\subset\C^2$ jollakin upotuksella $\C^2\subset\P^2$ määrittää RPI-pujoskaavion $\Omega$ $V$:n äärettömyyslinkille. Äärettömyyspisteiden lukumäärälle pätee $n_\infty(V)=n$, missä $n$ on kaavion juurisilmun valenssi. Lisäksi $\deg(V)=\ell_\text{root}$, jossa $\ell_\text{root}$ on juuren (virtuaalisen) komponentin linkkiluku äärettömyyslinkin kanssa. Jos $V$ on säännöllinen käyrä, $\Omega$ on säännöllinen (määritelmä???) pujoskaavio.

    (ii) Linkki on jonkin äärettömyyslinkin (joka voidaan valita säännölliseksi (määritelmä???), jopa hyväksi (määritelmä???)) osalinkki jos ja vain jos sillä on RPI-pujoskaavio.
\end{lause}
\begin{ajatus}
    Lause on sangen tiukasti muotoiltu, mutta sen ydinajatus on seuraava.
    \begin{enumerate}[i]
        \item Minkälainen solmu tai pujoskaavio algebrallisella käyrällä on?
        \item Millä ehdoilla algebralliset käyrät vastaavat linkkejä ja pujoskaavioita yksi yhteen?
    \end{enumerate}
\end{ajatus}